%% Chapter 13 ----------------------------------------------------------------
\chapter{Learmonth and STEREO/SWAVES Data}
\label{ch:other-radio}

Besides the e-CALLISTO archive, the application reads data from two further radio
instruments: the ground-based Learmonth Solar Radio Spectrograph in Western Australia, and
the space-based STEREO/WAVES (SWAVES) radio receivers. Both are opened from
\menu{Solar Events > Radio Bursts}.

\section{Learmonth Solar Radio Spectrograph}
\label{sec:learmonth}
\index{Learmonth}

The Learmonth Solar Radio Spectrograph, operated by the Australian Bureau of Meteorology,
records daily archive files (\file{.srs}\index{srs@.srs file}) in two receiver bands. \menu{Solar Events > Radio
Bursts > Learmonth} opens the \gui{Learmonth Solar Radio Browser}
(Figure~\ref{fig:learmonth}), which converts selected parts of a daily file into FIT files
for the main window.

\screenshot[0.7\linewidth]{radio-sources/learmonth_browser}{The Learmonth Solar Radio Browser.}{fig:learmonth}{The Learmonth Solar Radio Browser after Show Available Data, with the Daily Raw File information, the list of available 15-minute chunks (one highlighted for the chosen time) and the Convert buttons.}

\begin{description}
  \item[Observation Parameters] choose the \gui{Date} and a \gui{Time}, then click
    \gui{Show Available Data}. The daily file is read from the local cache when available
    and downloaded otherwise. The chunk containing the chosen time is highlighted.
  \item[Daily Raw File] information about the daily archive file. \gui{Download} saves
    the raw \file{.srs} file to a folder you choose.
  \item[Available 15-Minute Chunks] the time ranges into which the day is divided.
    \gui{Select All} and \gui{Deselect All} change the selection.
  \item[Convert to FIT] converts the selected chunks to CALLISTO-style FIT files in a folder
    you choose.
  \item[Convert and Import] converts the selected chunks and imports them into the main
    window, combined in time when consecutive.
\end{description}
Converted files carry a standard FITS header with the frequency and time axes and the
observatory position, so they can be processed and analyzed exactly like e-CALLISTO data.

\section{STEREO/SWAVES}
\label{sec:swaves}
\index{STEREO/SWAVES}\index{SWAVES}

SWAVES on the twin STEREO spacecraft \parencite{bougeret2008} records radio spectra from
2.6~kHz to 16~MHz, directly below the CALLISTO band. A burst that drifts out of the
ground-based range can therefore be followed into the interplanetary medium on the same
figure. \menu{Solar Events > Radio Bursts > SWAVES} opens the \gui{STEREO/SWAVES Radio
Spectrograph} dialog (Figure~\ref{fig:swaves-dialog}).

\screenshot[0.55\linewidth]{radio-sources/swaves_dialog}{The STEREO/SWAVES dialog.}{fig:swaves-dialog}{The STEREO/SWAVES Radio Spectrograph dialog with the Observation Window (UTC) fields, the Use CALLISTO Window, Cancel and Download and Plot buttons, and the summary line.}

\begin{description}
  \item[Start date, Start time, Duration] the UTC window to load (0.1--24~h; default
    4~h). The data are NASA/SPDF level-2 one-minute averages served as one file per UTC day;
    a window that crosses midnight fetches and joins both days. The archive starts on
    27~October~2006, and downloaded day files are cached and reused.
  \item[Spacecraft] \gui{STEREO-A (Ahead)}\index{STEREO} or \gui{STEREO-B (Behind)}. STEREO-B is
    unavailable for dates after 1~October~2014, when contact with the spacecraft was lost.
  \item[Sync padding] extra time added on each side when the window is taken from the
    CALLISTO spectrum (0--720~min; default 30~min). A Type~II or Type~III burst takes tens
    of minutes to hours to drift from the corona into the SWAVES band, so an exact match
    would cut off the part you want to see.
  \item[Use CALLISTO Window] copies the time range of the loaded CALLISTO spectrum, widened
    by the sync padding. \menu{Solar Events > Sync Current Time Window}\index{time synchronization} updates an open
    SWAVES dialog in the same way.
  \item[Download and Plot] loads the data; \gui{Cancel} stops a running load.
\end{description}

\subsection{The split view}
\label{sec:swaves-split}

Once SWAVES data are loaded, the plotting area of the main window splits into two panels:
CALLISTO above and SWAVES below, on a shared time axis, so that panning or zooming either
panel moves both (Figure~\ref{fig:swaves-split}). The CALLISTO interval is outlined on the
SWAVES panel. The SWAVES frequency axis is always logarithmic and its intensity is in
decibels above the instrument background, so the panel keeps its own color bar and color
scaling while following the color map chosen in the sidebar. Loading SWAVES with no
CALLISTO file open plots it across the whole area; it moves into the split view\index{STEREO/SWAVES!split view} as soon as a
FITS file is loaded.

\screenshot{radio-sources/swaves_split_view}{CALLISTO and STEREO/SWAVES spectra in the split view.}{fig:swaves-split}{Main window in split view: a CALLISTO spectrum with a Type II/III burst in the upper panel and the SWAVES spectrum below on the shared time axis, with the CALLISTO interval outlined.}

\menu{Solar Events > SWAVES Panel}\index{STEREO/SWAVES!panel} hides and restores the panel without discarding the
loaded data. The split view is included in exported figures, saved in and restored from
projects without downloading the data again, and added to project reports.

\begin{note}
  The drawing, lasso, drift and measurement tools act on the CALLISTO panel only, and the
  data-reduction controls of the sidebar do not apply to the SWAVES panel.
\end{note}
